
%==========================================================================

\begin{frame}[fragile]

  {\Huge Graphs}

  \vspace{10pt}

	{\large Structured Concurrency with Graphs.}

  \vspace{20pt}

  \textbf{Learning objectives:}
  \begin{itemize}
    \item {How to express dependencies between kernels with Graphs.}
    \item {The Graph phases: construction and use.}
    \item {Why Graphs can help with latency limited usecases.}
  \end{itemize}

  \vspace{-20pt}

\end{frame}

%==========================================================================

\begin{frame}[fragile]{What Are Graphs}
\textbf{Graphs as expression of Kernel dependencies}

\begin{itemize}
  \item{Ordering kernel execution so far: fence execution space instances in the right order.}
  \item{Graphs express this more directly}
  \item{Graphs are preconstructed, and then (repeatedly) submitted in their entirety for execution.}
  \item{Graphs can only be created on the host}
\end{itemize}

\pause

\textbf{Important Concepts:}

\begin{itemize}
  \item{Graph: the entire object}
  \item{Node: dependency points in the graph}
  \item{Root-Node: the starting point for the graph construction.}
\end{itemize}

%\pause
%\begin{block}{Important Point}
%  Execution Spaces execute operations in dispatch order.
%\end{block}

\end{frame}

%==========================================================================

\begin{frame}[fragile]{Creating Graphs}
\textbf{Creating a Graph via Scoped Setting}

\begin{itemize}
  \item{Scoping prevents mistakes}
	  \begin{itemize}
		  \item{Disallow modification of captured objects}
	  \end{itemize}
  \item{But less flexible when used in complex settings}
\end{itemize}

\pause
\begin{code}[linebackgroundcolor={},keywords={create_graph}]
auto graph = create_graph([&](const auto& root_node) {
  // build the graph
  ...
});
graph.submit();
\end{code}
\end{frame}
\begin{frame}[fragile]{Creating Graphs}
  \textbf{Creation via raw graph object}
\begin{itemize}
		  \item{More flexibility in terms of creating the graph in complex software}
		  \item{Need to be careful about capture semantics}
                  \item{Need to explicitly instantiate to be able to submit}
\end{itemize}

\begin{code}[linebackgroundcolor={},keywords={Graph}]
Kokkos::Experimental::Graph graph;
auto root_node = graph.root_node();

// build the graph
...

graph.instantiate();
graph.submit();
\end{code}


\end{frame}

%==========================================================================

\begin{frame}[fragile]{The Node Types}
\textbf{Attach nodes to each other - starting with root\_node}
\begin{itemize}
  \item{\texttt{then\_parallel\_for}: same as \texttt{parallel\_for}}
  \item{\texttt{then\_parallel\_reduce}: same as \texttt{parallel\_reduce}, but only \texttt{View} results}
  \item{\texttt{then}: attach a single work item}
  \item{\texttt{then\_host}: attach a single work item executing on host}
  \item{\texttt{when\_all}: a merge node - node depending on multiple previous ones}
\end{itemize}
\begin{code}[linebackgroundcolor={},keywords={}]
auto node_A = root_node.then_paralle_for("A", RangePolicy{0,N}, F_A);
auto node_B1 = node_A.then_paralle_reduce("B1", RangePolicy{0,N}, F_B, result_view_B);
auto node_B2 = node_A.then("B2",KOKKOS_LAMBDA() { ... });
auto node_B_merge = when_all(node_B1, node_B2);
auto node_C = node_B_merge.then_host("C", []() { printf("Host\n"); });
\end{code}
\end{frame}
%==========================================================================

\begin{frame}[fragile]{Graphs and Execution Spaces}
  \textbf{When not otherwise specified: DefaultExecutionSpace}
  \begin{itemize}
    \item{Specify Graph execution space type as template parameter}
    \item{Submit to specific instance in \texttt{submit}}
    \item{The graph is sequenced with other work in its submitted to instance}
    \item{Do NOT! pass execution space instances to policies for nodes.}
    \item{Supports multi-gpu too (ask for more details)}
  \end{itemize}

\begin{code}[linebackgroundcolor={},keywords={Graph,ExecutionSpace, exec_instance}]
Kokkos::Experimental::Graph<ExecutionSpace> graph;
auto root_node = graph.root_node();

// build the graph
...

graph.instantiate();
graph.submit(exec_instance);
\end{code}

\end{frame}
%==========================================================================

\begin{frame}[fragile]{Pitfalls}
  \textbf{No Execution During Constructions - Fixed Object After!}
  \begin{itemize}
    \item{During graph construction nodes are NOT executed}
    \item{The functors or lambdas are copied into graph during construction - and not updated later!}
  \end{itemize}
\vspace{-0.7cm}
\begin{code}[linebackgroundcolor={},keywords={}]

auto graph = create_graph([&](const auto& root_node) {
  Kokkos::View<int, SharedSpace> value("Value");
  value() = 1;
  auto node_A = root_node.then("A", KOKKOS_LAMBDA() {
    printf("A %i;\n",value()); value() += 3;
  });
  Kokkos::fence();
  auto scalar = value();
  printf("H %i;\n", scalar);
  return node_B = node_A.then("B", KOKKOS_LAMBDA() {
     printf("B %i %i;\n",scalar, value()); value() += 7;
    });
});
graph.submit();
graph.submit();
\end{code}

This prints: \texttt{H 1; A 1; B 1 4; A 11; B 1 14;}
\end{frame}

\begin{frame}[fragile]{Exercise: Execution Space Instances}
\textbf{Simulate MPI communication behavior with N neihbors}
\begin{code}
for(int n=0; n<neighbors; n++) {
  auto my_send_buffer = subview(my_send_buffer, n, ALL);
  auto my_recv_buffer = subview(my_recv_buffer, n, ALL);
  auto my_idxs = subview(idxs, n, ALL);
  pack_buffers(data, my_idxs, my_send_buffer);
  exchange_messages(my_send_buffer, my_recv_buffer);
  unpack_buffers(data, my_recv_buffer);
}
\end{code}


  \textbf{Details}:
  \begin{small}
  \begin{itemize}
\item Location: \ExerciseDirectory{graph}
\item Create a graph
\item Pass nodes into the functions
\item Use a host node to print a message during exchange
\end{itemize}
  \end{small}

\ul{\textbf{Things to try:}}
  \begin{small}
  \begin{itemize}
  \item Vary problem size, number of neighors and instances (-S; -N; -I)
  \end{itemize}
  \end{small}

\end{frame}


%==========================================================================

